\begin{enumerate}
	\item Sur l'intervalle $\left [0~;~2 \pi\strut \right ]$, l'équation  $\sin (x) = 0,1$ admet :

	\begin{tabularx}{\linewidth}{*{2}{>{\bfseries\arraybackslash}p{5mm}@{}X}}
		a.& zéro solution&
		b.& une solution\\
		c.& deux solutions&
		d.& quatre solutions\\
	\end{tabularx}

On peut par exemple dire que la courbe représentant la fonction sinus et la droite d'équation $y=0,1$ ont deux points d'intersection sur l'intervalle $\left [ 0\;;\; 2\pi\strut \right ]$.

\begin{center}
\psset{xunit=1cm,yunit=1cm}
\begin{pspicture}(-1,-1.5)(7,1.5)
\psaxes[trigLabels=true,trigLabelBase=2,dx=\psPiH,xunit=\psPi]{->}(0,0)(0,-1.5)(2.2,1.5)
\psplot[linewidth=1.2pt,linecolor=red,plotpoints=5000]{0}{6.2832}{x RadtoDeg sin}
\psline[linecolor=blue](0,0.1)(6.2832,0.1)
\end{pspicture}
\end{center}

\hfill\textbf{Réponse c.}

	\item On considère la fonction $f$ définie sur l'intervalle $\left [0~;~\pi\strut\right ]$ par  $f(x) = x + \sin (x).$


	\begin{tabularx}{\linewidth}{{>{\bfseries\arraybackslash}p{5mm}@{}X}}
		a.& La fonction $f$ est convexe sur l'intervalle $[0~;~\pi]$\\
		b.& La fonction $f$ est concave sur l'intervalle $[0~;~\pi]$\\
		c.& La fonction $f$ admet sur l'intervalle $[0~;~\pi]$ un unique point d'inflexion\\
		d.& La fonction $f$ admet sur l'intervalle $[0~;~\pi]$ exactement deux points d'inflexion\\
	\end{tabularx}

$f'(x)=1+\cos(x)$ et $f''(x)=-\sin(x)\leqslant 0$ sur $\left [0~;~\pi\strut\right ]$, donc la fonction $f$ est concave sur cet intervalle.

\hfill\textbf{Réponse b.}

	\item Une urne contient cinquante boules numérotées de 1 à 50. On tire successivement trois boules dans cette urne, \textbf{sans remise}. On appelle \og{} tirage \fg{} la liste non ordonnée des numéros des trois boules tirées.
Quel est le nombre de tirages possibles, \textbf{sans tenir compte de l'ordre des numéros}?

	\begin{tabularx}{\linewidth}{*{4}{>{\bfseries\arraybackslash}p{5mm}@{}X}}
		a.& $50^{3}$&
		b.& $1 \times 2 \times 3$&
		c.& $50 \times 49 \times 48$&
		d.& $\dfrac{50 \times 49 \times 48}{1 \times 2 \times 3}$\\
	\end{tabularx}

$\ds\binom{50}{3}=\dfrac{50\times 49 \times 48}{1 \times 2 \times 3}$
\hfill\textbf{Réponse d.}

	\item On effectue dix lancers d'une pièce de monnaie. Le résultat d'un lancer est \og{}pile\fg{} ou \og{}face\fg{}. On note la liste ordonnée des dix résultats.\\
Quel est le nombre de listes ordonnées possibles?

	\begin{tabularx}{\linewidth}{*{2}{>{\bfseries\arraybackslash}p{5mm}@{}X}}
		a.& $2 \times 10$&
		b.& $2^{10}$\\
		c.& $1 \times 2 \times 3 \times \cdots \times 10$&
		d.& $\dfrac{1 \times 2 \times 3 \times \cdots \times 10}{1 \times 2}$\\
	\end{tabularx}

Il y a 2 résultats possibles si on effectue 1 lancer, $2^2$ résultats possibles si on effectue 2 lancers, etc., $2^{10}$ résultats possibles si on effectue 10 lancers.

\hfill\textbf{Réponse b.}

	\item On effectue $n$ lancers d'une pièce de monnaie équilibrée. Le résultat d'un lancer est \og{}pile\fg{} ou \og{}face\fg{}. On considère la liste ordonnée des $n$ résultats.\\
Quelle est la probabilité d'obtenir au plus deux fois \og{} pile \fg{} dans cette liste ?

	\begin{tabularx}{\linewidth}{*{2}{>{\bfseries\arraybackslash}p{5mm}@{}X}}
		a.& $\dfrac{n(n-1)}{2}$&
		b.& $\dfrac{n(n-1)}{2} \times\left(\dfrac{1}{2}\right)^{n}$\\
		c.& $1 + n + \dfrac{n(n-1)}{2}$&
		d.& $\left(1 + n + \dfrac{n(n-1)}{2}\right) \times\left(\dfrac{1}{2}\right)^{n}$
	\end{tabularx}
	
La variable aléatoire $X$ qui donne le nombre de résultats \og pile \fg{} sur $n$ lancers suit la loi binomiale de paramètres $n$ et $\dfrac{1}{2}$.

On cherche $P(X\leqslant 2)$ soit $P(X=0)+P(X=1)+P(X=2)$.
 
\textbullet~~$P(X=0)=\ds\binom{n}{0}\left (\dfrac{1}{2}\right )^0 \left (1-\dfrac{1}{2}\right )^{n-0} = \left (\dfrac{1}{2}\right )^n$

\textbullet~~$P(X=1)=\ds\binom{n}{1}\left (\dfrac{1}{2}\right )^1 \left (1-\dfrac{1}{2}\right )^{n-1} = n \left (\dfrac{1}{2}\right )^n$

\textbullet~~$P(X=2)=\ds\binom{n}{2}\left (\dfrac{1}{2}\right )^2 \left (1-\dfrac{1}{2}\right )^{n-2} = \dfrac{n(n-1)}{2} \left (\dfrac{1}{2}\right )^n$

Donc $P(X\leqslant 2)= \left (\dfrac{1}{2}\right )^n + n \left (\dfrac{1}{2}\right )^n + \dfrac{n(n-1)}{2}\left (\dfrac{1}{2}\right )^n = \left (1+n+\dfrac{n(n-1)}{2}\right )\left (\dfrac{1}{2}\right )^n$.

\hfill\textbf{Réponse d.}
\end{enumerate}
